\documentclass[12pt]{article}
\usepackage[utf8]{inputenc}
\usepackage[T1]{fontenc}
\usepackage[margin=1in]{geometry}
\usepackage{setspace}
\usepackage{xcolor}
\usepackage[colorlinks=true, urlcolor=blue!60!black, linkcolor=blue!60!black]{hyperref}
\setlength{\parindent}{0em}
\setlength{\parskip}{0.45em}
\renewcommand{\baselinestretch}{1.0}
\begin{document}
\begin{center}
{\Large \textit{United States v. Morrison}, 529 U.S. 598 (2000)\par}
\vspace{0.3em}
{\normalsize Case Brief --- One-Page Edition\par}
\end{center}

\textbf{Memory Jogger.} VAWA's civil remedy fails: noneconomic violence cannot be aggregated, and Section 5 does not reach private actors (\href{https://www.courtlistener.com/opinion/118363/united-states-v-morrison/}{CourtListener}).

\textbf{Facts.} Virginia Tech student Christy Brzonkala alleged that football players Antonio Morrison and James Crawford raped her; the university reduced and then set aside Morrison's suspension. She sued under 42 U.S.C. \S~13981, VAWA's federal damages action for ``crimes of violence motivated by gender,'' which rested on four years of hearings.

\textbf{Procedural History.} The United States intervened to defend the law. The district court held \S~13981 unconstitutional; a Fourth Circuit panel reversed, but the en banc court affirmed.

\textbf{Votes.} Affirmed 5--4. Rehnquist wrote for the Court; Thomas concurred; Souter and Breyer dissented, with Stevens and Ginsburg.

\textbf{Holding.} Congress may not regulate noneconomic violent conduct based on aggregate effects, even with extensive findings; Section 5 reaches only state action, and a nationwide remedy against private individuals is not congruent and proportional.

\textbf{Reasoning.} Gender-motivated violence is ``not, in any sense of the phrase, economic activity,'' and the findings' but-for chain would reach all crime and family law. The remedy burdened no state official and applied even where no state bias was shown. Thomas would scrap the substantial-effects test. Souter cited Congress's ``mountain of data'' and pre-1937 formalism; Breyer called the economic line unworkable.

\textbf{Examples.} \emph{Same-Side:} a federal tort for local hate crimes. \emph{Opposite-Side:} interstate domestic-violence crimes. \emph{Fence-Sitter:} an intrastate endangered species.

\textbf{Critique.} Randy Barnett praised enforcement of enumerated limits; Robert Post and Reva Siegel said the Court treated gender violence as ``local'' while monopolizing Fourteenth Amendment enforcement.

\textbf{Key Quotation.} ``We accordingly reject the argument that Congress may regulate noneconomic, violent criminal conduct based solely on that conduct's aggregate effect on interstate commerce. The Constitution requires a distinction between what is truly national and what is truly local.'' --- Chief Justice Rehnquist, 529 U.S. 598.

\end{document}
